% This ensures VS Code compiles the main file even when you're editing a chapter.
% !TeX root = ../main.tex

\chapter{Proposal}
\label{cap:prop}

This chapter presents the proposed architecture of the prototype used to test the
feasibility of a sensor fusion system for precise ground target localization. As mentioned
in Chapter~\ref{cap:intro}, this work focuses on the development of the drone pod subsystem,
specifically on the sensor fusion for precise distance measurements and parallax correction.

As the drone itself is not as critical as the payload, it will be replaced by an
embedded computer that will execute the functions of the drone. It will not only perform
computer vision tasks and command the distance-measuring pod, but also receive all the data
processed by the pod and log everything, allowing the system's performance to be analyzed.

As mentioned in Chapter~\ref{cap:intro}, the most commonly used electronic sensors for
distance measurements in drones are single-point or multi-point \gls{tof} sensors. As the
pod is to intended to be used as a standalone, and as flexible as possible as a drone
component, it must be able to operate over any distance that a drone might need to measure
for a ground object.
Assuming that the drone will fly at an approximate maximum altitude of 120m, this means that
the target will be at least 120m away from the drone, under ideal conditions,
when it is directly below the sensor.
As the target moves away from the point directly below the drone on the ground,
the distance increases. This relationship between the position of the ground target and the
distance to the drone can be better understood via the Figure~\ref{fig:3img1}. The final
distance can be expressed either by the Equation~\ref{eq:dist_to_gnd1} or by the
Equation~\ref{eq:dist_to_gnd2}:

\begin{figure}[ht]
  \centering
  \includegraphics[width=0.55\textwidth]{4/img1}
  \caption{Distance from drone to target as a function of target's position on ground in 2D}
  \label{fig:3img1}
\end{figure}

\begin{equation}
  \label{eq:dist_to_gnd1}
  \scalebox{1.5}{$ d = \sqrt{ x^2 + h^2 } $}
\end{equation}

\begin{equation}
  \label{eq:dist_to_gnd2}
  \scalebox{1.5}{$ d = h \cdot sec\bigg( ~ tan^-1\Big( ~ \frac{x}{h} ~ \Big) ~ \bigg) $}
\end{equation}

%where:

\begin{itemize}
  \item $d$ is the distance to the object on the ground
  \item $h$ is the height of the drone
  \item $x$ is the position of the object on the ground
\end{itemize}

With such large distances like that, flash \gls{lidar} becomes too expensive or
difficult to source.
As the pod will need to provide only the distance to a single target at a time,
there is no need to use a complex system such as scanning \gls{lidar}, as it is it's
heavy, consumes a significant amount of power, and has a continuously running motor.
As mentioned in the Chapter~\ref{cap:intro}, a hybrid solution using a single \gls{tof}
sensor mounted on a gimbal achieves both the ability to point in multiple directions
and to operate over long distances with cheaper and simpler devices.

\section{Drone subsystem}

As already mentioned in this chapter as well as in Chapter~\ref{cap:back}, the drone
subsystem will be responsible for behaving as a drone, so the performance of the distance
measuring pod can be evaluated indoors, without the need to cycle batteries:

\begin{itemize}
  \item Image acquisition from the onboard camera.
  \item Image processing and feature detection.
  \item Target localization.
  \item Gimbal pointing control.
  \item Distance information logging.
\end{itemize}

The drone subsystem will consist of a \gls{sbc}, connected to the camera and the pod,
running a program that executes the computer vision task of identifying the target and
its respective direction from the drone.

\subsection{SBC and Camera}

The chosen \gls{sbc} for the project was the Raspberry Pi model 5 8GB \cite{rasppi51}.
This \gls{sbc} was chosen because it is the latest iteration of the Raspberry Pi \gls{sbc}
family, which holds an approximate 30\% share of the \gls{sbc} market \cite{rasppi52}.

As the Raspberry Pis are the most popular \gls{sbc} in the market, it is the
easiest \gls{sbc} to purchase anywhere in the world.
Also, because of their popularity, there is plenty of information and support online,
not only in official forums but also on other engineering forums and sources.
This drastically facilitates development.

This \gls{sbc} includes a
Broadcom BCM2712 2.4GHz quad-core 64-bit Arm Cortex-A76 CPU,
VideoCore VII GPU,
8GB LPDDR4X-4267 SDRAM,
and many interfaces and connectors, including 4-lane MIPI connectors for the camera and
a pin header for many different serial protocols to be used to communicate with the pod
subsystem. The board is powered by a 5V/5A DC power supply via USB-C.
The exact model can be seen in Figure~\ref{fig:4rasp}.

\begin{figure}[h]
  \centering
  \includegraphics[width=0.5\textwidth]{4/rasp}
  \caption{Raspberry Pi model 5}
  \label{fig:4rasp}
\end{figure}

The \gls{sbc} will run on Raspberry Pi OS, an OS made specifically for the Raspberry Pi
hardware. Raspberry Pi OS, known as Raspbian, is a port of Debian Trixie. Debian
itself is a Linux distribuition, which makes Raspbian also a Linux distribuition.

The chosen camera to be connected to the Raspberry Pi 5 was the Camera Module 3, an
official camera module for the Raspberry Pi \gls{sbc} family \cite{raspcam}. The camera
uses a back-illuminated, stacked CMOS 12-megapixel Sony IMX708 image sensor, capable
of capturing images with 4608 $\times$ 2592 pixels, with a diagonal field of view of 75
degrees. The chosen resolution used will be 2304 $\times$ 1296 pixels, so it is still
very close to FullHD resolution but smaller, which may help with processing time when
running computer vision algorithms, while keeping the same ratio, thus the same
\gls{fov}.

\subsection{Computer vision subroutine}

For the computer vision task, the library chosen was \gls{opencv}, in Python.
OpenCV is the most widely used computer vision library in the world \cite{opencv}.
It has functions to perform every part of the computer vision pipeline involved in
the project, from image acquisition to feature detection.
This library will be used to obtain simple feature detection capabilities,
allowing the project to be tested with tasks such identifying a dot on the
wall, or a circle of certain color.
This will make testing the whole system a lot easier by being able to print targets
on paper and fixed them on many different surfaces.

The chosen programming language to be used in the drone subsystem is Python, due to
it ease of learning, use, and test many different algorithms, including \gls{opencv}.
The Python code will make calls to Raspbian to receive images from the camera, and to
send messages to the pod subsystem via the communication pin header.

\section{Distance measuring pod}

The second subsystem in the prototype, and the main focus of this work, the distance
measuring pod will consist mainly of a distance measuring sensor, a gimbal to aim the
sensor according to a specific direction, a microcontroller to control everything,
process the data and send it back to the drone subsystem.
The full cycle of the measuring routine would consist of:

\begin{itemize}
  \item Direction data communication from the drone.
  \item Gimbal control.
  \item Sensor data acquisition.
  \item Data processing and parallax compensation.
  \item Final distance data communication back to the drone.
\end{itemize}

In order to perform all these tasks, a digital, programmable system will be used. The
hardware chosen for this task is a microcontroller. A microcontroller is a compact
integrated circuit designed to perform specific control tasks within an embedded system.
It combines a processor core, memory (such as RAM and flash), and input/output
peripherals on a single chip.
Unlike general-purpose computers, microcontrollers are optimized for
interacting with hardware, allowing them to read sensors, process data, and control
actuators in real time.

\subsection{Microcontroller and framework}

The chosen microcontroller for the pod subsystem is the ESP32-S3 \cite{esp32s3}.
The ESP32-S3 is a low-power microcontroller developed by Espressif Systems.
It is the successor of the well known ESP32. The ESP32-S3 features
a dual-core XTensa LX7 processor architecture, running at 240MHz,
512kB of SRAM,
and 384 KB ROM.
The device features a wide range of peripherals, such as GPIO, SPI, I²C, UART, ADC,
and PWM, enabling flexible interfacing with sensors and external components.
A board with the chosen microcontroller module can be seen in Figure~\ref{fig:4esp}.

\begin{figure}[h]
  \centering
  \includegraphics[width=0.5\textwidth]{4/esp}
  \caption{ESP32-S3 module on development board}
  \label{fig:4esp}
\end{figure}

The microcontroller was chosen for its availability, flexibility, price point and also
the development environment. Espressif Systems developed a framework called ESP-IDF
\cite{espidf}. This framework includes a \gls{hal}, which gives the embedded developer
a fast and reliable way of accessing the hardware controlled by the microcontroller.
This shortens development time and system bring-up tests, which helps with the rapid
prototyping phase of the project.

\subsection{Single-point ToF sensor}

The chosen sensor to measure distance to objects from the pod is a ToFFuture Technology
XT-S1 \gls{lidar} sensor \cite{xts1}. The XT-S1 is a compact single-point \gls{ir} laser sensor
based on the \gls{tof} principle.
It weighs only 9 grams and consumes less than 0.5W.
The sensor can operate in environments with intense lighting conditions of up to
100klux, while measuring distance up to 24m distant objects of 10\% reflectivity.
It's accuracy is $\pm$2\%, with a measurement frequency of 50Hz.

The sensor was chosen for its easy availability, low price, small size and weight and
good outdoor performance. The sensor module can be seen in Figure~\ref{fig:4xts1}.

\begin{figure}[h]
  \centering
  \includegraphics[width=0.5\textwidth]{4/xts1}
  \caption{Xt-S1 sensor module and it's dimensions}
  \label{fig:4xts1}
\end{figure} 

\subsection{Sensor gimbal}

A gimbal is a mechanical system that allows a sensor to rotate about one or more axes,
typically to stabilize its orientation or to actively control its field of view.
By decoupling the sensor from the motion of the base platform, a gimbal enables precise
pointing and smooth tracking, which is essential in applications such as imaging,
scanning, and navigation.

The motion of a gimbal is commonly achieved using servomotors, which provide controlled
angular positioning through closed-loop feedback. In a typical configuration, each axis
of the gimbal is driven by a servomotor that receives a position command and adjusts its
shaft angle accordingly. Feedback is obtained either from an internal sensor or an
external encoder, allowing the control system to minimize the error between the desired
and actual position. This enables accurate, repeatable, and dynamic movement of
the sensor.

Traditional servomotors are often based on brushed DC motors combined with a gearbox and
a position feedback element, such as a potentiometer. These systems are relatively simple
and cost-effective, but they suffer from mechanical wear due to the brushes in the long
term, limited precision from the feedback device, and backlash introduced by the gearbox.

In contrast, modern servomotors based on brushless DC (BLDC) motors use electronic
commutation and typically incorporate high-resolution magnetic angle sensors for position
feedback. These sensors provide precise, contactless measurement of the motor's shaft
angle, eliminating wear and improving reliability. BLDC-based servos also offer smoother
motion, higher efficiency, lower noise, and better dynamic response compared to brushed
counterparts. The downsides of using BLDC motors in a servo system are a significant
increase in complexity, price, and availability of components.

The chosen type of motor for the prototype will initially be brushed DC servomotors, as
they are cheaper, easy to purchase, and easy to model and control, in contrast with their
BLDC counterparts, which usually need customized control solutions and BLDC drivers. In
order to compensate for the downsides of using brushed DC servomotors, high-torque,
robust servomotor models were used, for these have smaller mechanical tolerances, show
less wear and tear due to improved quality on gears and more accurate digital feedback
sensor systems.

\subsection{Parallax compensation}

% fuck me i guess

\section{3D printing}

Additive manufacturing, specifically 3D printing, is employed in the development of the
prototype chassis for both the drone and pod subsystems as a means of rapid prototyping
and design iteration. This approach enables the quick fabrication of structural
components directly from digital models, significantly reducing development time and
cost compared to traditional manufacturing methods.

Through 3D printing, it is possible to produce custom-designed parts that precisely
match the geometry and requirements of the system. This is particularly advantageous for
integrating sensors, electronics, and mechanical assemblies, as mounting features,
enclosures, and supports can be tailored to fit specific components without the need for
standardized parts or complex machining processes.

Additionally, the use of 3D printing allows for iterative testing and refinement of the
chassis design. Components can be easily modified and reprinted to address issues such
as structural integrity, weight distribution, or accessibility. This flexibility is
especially valuable during the early stages of development, where design requirements
may evolve rapidly. 
